% =============================================================================
% Kapitel 8: Diskussion
% =============================================================================

Dieses Kapitel ordnet die Evaluationsergebnisse aus Kap.~\ref{sec:evaluation} in den Kontext der Forschungsfrage ein, diskutiert die Implikationen des hybriden Ansatzes und benennt die Limitationen der Arbeit.

\subsection{Beantwortung der Forschungsfrage}\label{sec:beantwortung}

Die Forschungsfrage lautet: \textit{Welche Potenziale und Grenzen bietet ein Embedding-basierter Ansatz gegenüber regelbasierter Textanalyse für die Unterstützung Lehrender bei der Erstellung qualitativ hochwertiger Kompetenzformulierungen in deutschsprachigen Modulhandbüchern?}

Die Evaluation auf dem Gold-Standard-Datensatz (5\,617~Sätze aus Modulhandbüchern mehrerer Hochschulen, Schwerpunkt Hochschule Fulda) erlaubt eine differenzierte Antwort entlang der vier Analysedimensionen.

\paragraph{Potenziale.}
Das größte Potenzial zeigt sich bei der \textbf{Satztyp-Klassifikation}: Der Embedding-Ansatz erreicht $F_1 = 0{,}963$ gegenüber $F_1 = 0{,}819$ beim regelbasierten Ansatz.
Besonders ausgeprägt ist der Vorteil bei Sätzen, die der regelbasierte Ansatz strukturell nicht verarbeiten kann -- Inhaltsbeschreibungen ($F_1 = 0{,}645$ vs.\ $0{,}000$) und Sonstiges ($F_1 = 0{,}778$ vs.\ $0{,}177$).
Der regelbasierte Ansatz kann ohne Verbliste-Evidenz keine Negativklassifikation vornehmen; die semantische Repräsentation durch Sentence-Embeddings ermöglicht diese Unterscheidung erstmals.

Bei der \textbf{Taxonomie-Zuordnung} übertrifft der Embedding-Ansatz ($F_1 = 0{,}839$) den regelbasierten ($F_1 = 0{,}581$) ebenfalls deutlich.
Entscheidend ist nicht nur die höhere Genauigkeit, sondern die vollständige Abdeckung: Der regelbasierte Ansatz kann nur 48{,}2\,\% der K-Sätze überhaupt zuordnen, da er auf exakte Verbliste-Treffer angewiesen ist.
Der Embedding-Ansatz klassifiziert alle Sätze und erreicht dabei in jeder der sechs Taxonomiestufen einen höheren F1-Wert.
Stufe~1 (Erinnern) illustriert dies besonders deutlich: $F_1 = 0{,}860$ vs.\ $0{,}129$.

Das \textbf{Empfehlungssystem} stellt eine qualitativ neue Funktion dar, die mit dem regelbasierten Ansatz prinzipbedingt nicht realisierbar ist.
Mit Precision@1 $= 0{,}510$ liefert die Hälfte der Empfehlungen eine taxonomisch passende Referenzformulierung.

Die \textbf{Set-Analyse} profitiert direkt von der verbesserten Einzelsatz-Klassifikation: Nur mit vollständiger Abdeckung kann die Taxonomie-Verteilung eines Moduls zuverlässig bestimmt werden.

\paragraph{Grenzen.}
Die Grenzen des Embedding-Ansatzes zeigen sich in drei Bereichen.
Erstens bleibt Stufe~4 (Analysieren) für alle Ansätze die schwierigste ($F_1 = 0{,}667$ beim Embedding-Ansatz), was auf die semantische Nähe zu Stufe~3 (Anwenden) und die geringe Trainingsdatenmenge zurückzuführen ist.
Zweitens ist der Embedding-Ansatz nicht domänenunabhängig: Bei TH Mittelhessen (Ingenieurwesen) sinkt die Taxonomie-Zuordnung auf $F_1 = 0{,}717$, was die von Huber und Niklaus~\cite{huber2025llmsblooms} berichtete Domänenabhängigkeit von Embedding-Klassifikatoren bestätigt.
Drittens fehlt dem Embedding-Ansatz die Transparenz des regelbasierten: Die Zuordnung eines Verbs zu einer Taxonomiestufe über eine statische Liste ist für Lehrende unmittelbar nachvollziehbar, eine SVM-Entscheidung auf einem 768-dimensionalen Embedding nicht.

\subsection{Kosten-Nutzen des hybriden Ansatzes}\label{sec:kosten-nutzen}

Der hybride Ansatz wurde mit der Hypothese konzipiert, die Stärken beider Verfahren zu kombinieren: die Transparenz und Geschwindigkeit des regelbasierten Ansatzes für eindeutige Fälle, die Kontextsensitivität der Embeddings für unsichere Fälle.
Die Evaluation zeigt, dass diese Hypothese nur teilweise zutrifft.

\paragraph{Accuracy.}
In beiden Hauptdimensionen liegt der hybride Ansatz zwischen den anderen beiden: Typ-Klassifikation $F_1 = 0{,}929$ (vs.\ Embedding $0{,}963$), Taxonomie $F_1 = 0{,}713$ (vs.\ Embedding $0{,}839$).
Der Abstand zum reinen Embedding-Ansatz beträgt 3{,}4 bzw.\ 12{,}6~Prozentpunkte.
Ursache ist die Fehlerfortpflanzung im Cascading: Wenn die regelbasierte Stufe einen Satz fälschlich als K klassifiziert, wird das Embedding-Fallback nicht aktiviert.
Die Quellen-Verteilung zeigt, dass 58{,}2\,\% der Sätze im Test-Set regelbasiert klassifiziert werden -- darunter auch fehlerhafte Zuordnungen, die das Embedding hätte korrigieren können.

\paragraph{Laufzeit.}
Der potenzielle Laufzeitvorteil des hybriden Ansatzes ergibt sich daraus, dass für einen Teil der Sätze keine Embedding-Berechnung erforderlich ist.
In der Praxis ist dieser Vorteil gering: Die Embedding-Berechnung mit dem SBERT-Modell benötigt auf moderner Hardware wenige Millisekunden pro Satz, und der Flaschenhals liegt bei der initialen Modellladung (ca.\ 3~Sekunden), nicht bei der Inferenz.
Da das Empfehlungssystem für jeden K-Satz ohnehin eine Ähnlichkeitssuche auf Embedding-Basis durchführt, wird das SBERT-Modell in jedem Fall geladen und die Embeddings aller Sätze berechnet.
Der Laufzeitvorteil des Cascading reduziert sich damit auf die eingesparte SVM-Inferenz für regelbasiert zugeordnete Sätze -- eine Operation im Mikrosekundenbereich.

\paragraph{Fazit.}
Der hybride Ansatz bringt in der aktuellen Cascading-Architektur keinen messbaren Vorteil gegenüber dem reinen Embedding-Ansatz: Er ist weniger genau und nicht signifikant schneller.
Sein Mehrwert liegt in der Transparenz: Für Sätze mit eindeutigen Verbliste-Treffern kann das System die Zuordnung erklären (\enquote{Das Verb X ist der Stufe Y zugeordnet}), was für die Akzeptanz durch Lehrende relevant sein kann.

\subsection{Cascading vs.\ alternative Hybridisierungsstrategien}\label{sec:cascading-alternativen}

Das gewählte Cascading-Muster ist nicht die einzige Möglichkeit, regelbasierte und Embedding-basierte Analyse zu kombinieren.
Zwei Alternativen verdienen eine Diskussion.

\paragraph{Oder-Modell (Union).}
Beide Ansätze werden parallel ausgeführt, und ein Satz wird als K klassifiziert, wenn mindestens einer der beiden Ansätze ihn als K einstuft.
Dies maximiert den Recall, kann aber zu einer Reduktion der Precision führen, da Falsch-Positive beider Ansätze kumuliert werden.

\paragraph{Und-Modell (Intersection).}
Ein Satz wird nur als K klassifiziert, wenn beide Ansätze übereinstimmen.
Dies maximiert die Precision auf Kosten des Recalls: Sätze, die nur einer der beiden Ansätze erkennt, gehen verloren.

\paragraph{Voting / Ensemble.}
Beide Ansätze geben ihre Klassifikation ab; bei Übereinstimmung wird das Ergebnis übernommen, bei Dissens entscheidet eine gewichtete Kombination (z.\,B.\ basierend auf der SVM-Konfidenz).
Dieser Ansatz adressiert die Fehlerfortpflanzung des Cascading, erfordert aber stets die Ausführung beider Ansätze und eliminiert damit den theoretischen Laufzeitvorteil.

Da der reine Embedding-Ansatz bereits $F_1 = 0{,}963$ (Typ) bzw.\ $0{,}839$ (Taxonomie) erreicht, ist der erwartbare Zugewinn durch ein Ensemble gering.
Die Ergebnisse legen nahe, dass die regelbasierte Komponente ihren Mehrwert primär in der Erklärbarkeit hat, nicht in der Klassifikationsgenauigkeit.
Ein vielversprechender Ansatz für zukünftige Arbeiten wäre, die regelbasierte Zuordnung als \emph{Erklärungsschicht} zu verwenden: Das Embedding klassifiziert, und wenn ein Verbliste-Treffer vorliegt, wird die Zuordnung mit der Listenregel erklärt.

\subsection{Einordnung in den Forschungskontext}\label{sec:einordnung}

Die erzielten Ergebnisse sind im Kontext verwandter Arbeiten plausibel.
Li et al.~\cite{li2022bloomclassification} berichten F1-Werte zwischen 0{,}88 und 0{,}95 für die Bloom-Klassifikation auf einem englischsprachigen Datensatz von 21\,380~Lernergebnissen -- der hier erzielte $F_1 = 0{,}839$ für die Taxonomie-Zuordnung auf dem deutschen Datensatz (5\,174~K-Sätze) liegt im gleichen Bereich, obwohl die deutsche Sprache und die geringere Datensatzgröße zusätzliche Herausforderungen darstellen.
Kumar et al.~\cite{kumar2025automated} bestätigen mit 94\,\% Accuracy für SVM auf Sentence-Embeddings die Wirksamkeit dieses Ansatzes gerade auf kleineren Datensätzen.

Der von Huber und Niklaus~\cite{huber2025llmsblooms} beobachtete Domänentransfer-Verlust zeigt sich auch in den hier vorliegenden Daten: Die Taxonomie-Zuordnung bei TH Mittelhessen ($F_1 = 0{,}717$) liegt deutlich unter dem Wert für Hochschule Fulda ($F_1 = 0{,}911$).
Der Datensatz dieser Arbeit umfasst mit vier Fachbereichen eine größere Domänenbreite als die genannten Arbeiten, wobei die Generalisierbarkeit durch die Konzentration auf eine Hauptinstitution (57\,\% Fulda) eingeschränkt bleibt.

\subsection{Limitationen}\label{sec:limitationen}

Die Ergebnisse unterliegen mehreren Einschränkungen.

\paragraph{Datensatz.}
Der Gold-Standard-Datensatz ist mit 92{,}1\,\% Kompetenzformulierungen stark unbalanciert.
Die Minderheitsklassen I ($n = 188$) und S ($n = 255$) sind unterrepräsentiert, was die F1-Werte für diese Klassen mit Vorsicht interpretieren lässt.
Die Stichproben externer Hochschulen sind klein (TU Darmstadt: 180, Uni Kassel: 7~im Test-Set), sodass die Generalisierbarkeitsaussagen eingeschränkt sind.

\paragraph{Annotation.}
Die LLM-gestützte Vorannotation birgt das Risiko systematischer Verzerrungen.
Obwohl die Inter-Annotator-Reliabilität für Satztyp ($\kappa = 0{,}737$) und Taxonomie ($\kappa = 0{,}861$) substanziell bis fast perfekt ist, zeigt die Qualitätsdimension ($\kappa = 0{,}350$/$0{,}401$ vs.\ LLM) eine deutliche Diskrepanz zwischen menschlicher und maschineller Bewertung.
Da der gesamte Datensatz mit LLM-Labels trainiert und nur 300~Sätze manuell validiert wurden, könnten systematische LLM-Fehler in den nicht-validierten 5\,317~Sätzen die Trainingsdaten verzerren.

\paragraph{Qualitätsdimension.}
Die Formulierungsqualität wird in dieser Arbeit zwar annotiert, aber nicht als automatische Analysefunktion evaluiert.
Die geringe Übereinstimmung zwischen LLM und menschlichen Annotatoren ($\kappa = 0{,}350$ bzw.\ $0{,}401$) deutet darauf hin, dass die dreistufige Skala (gut/akzeptabel/schlecht) ohne operationalisierte Kriterien zu subjektiv ist.

Der HRK-nexus-Leitfaden~\cite{hrknexus2015lernergebnisse} definiert sieben Leitlinien für die Formulierung von Lernergebnissen, aus denen sich messbare Qualitätskriterien ableiten lassen:

\begin{enumerate}
  \item \textbf{Beobachtbares Verb}: Enthält der Satz genau ein beobachtbares Handlungsverb aus einer Taxonomie?
  \item \textbf{Eindeutigkeit}: Ist das Verb eindeutig einer Taxonomiestufe zuordenbar (keine Mehrdeutigkeit)?
  \item \textbf{Subjekt-Nennung}: Werden die Studierenden als Subjekt benannt oder impliziert?
  \item \textbf{Spezifität}: Enthält der Satz einen konkreten fachlichen Gegenstand (nicht nur das Verb)?
  \item \textbf{Messbarkeit}: Ist das beschriebene Ergebnis feststell- und überprüfbar?
  \item \textbf{Granularität}: Wird pro Satz nur ein Lernergebnis formuliert?
  \item \textbf{Taxonomie-Abdeckung}: Sind auf Modulebene verschiedene kognitive Stufen vertreten?
\end{enumerate}

Die Kriterien~1--4 und~6 lassen sich teilweise automatisiert prüfen (Verbliste-Treffer, Anzahl Verben, Subjekt-Erkennung per spaCy, Satzlänge als Proxy für Spezifität).
Kriterium~5 (Messbarkeit) erfordert eine inhaltliche Beurteilung, für die Embedding-basierte Ansätze einen Ausgangspunkt bieten könnten.
Kriterium~7 wird bereits durch die Set-Analyse adressiert.
Eine Reannotation des Datensatzes anhand dieser operationalisierten Kriterien -- statt der pauschalen dreistufigen Skala -- würde eine reliable Qualitätsbewertung ermöglichen und ist als nächster Schritt vorgesehen.

\paragraph{Semantische Affinität zwischen Annotation und Evaluation.}
Ein grundsätzliches methodisches Risiko ergibt sich aus der Verwendung LLM-generierter Labels sowohl für das Training als auch für die Evaluation.
Da 94{,}7\,\% des Datensatzes ausschließlich LLM-annotiert sind, testet der Vergleich der Ansätze teilweise, wie gut ein semantisches Verfahren (SBERT+SVM) die Urteile eines anderen semantischen Verfahrens (Claude Sonnet) reproduziert.
Der regelbasierte Ansatz funktioniert kategorial anders (exakter Stringvergleich) und könnte durch diese Konstellation systematisch benachteiligt werden -- nicht weil seine Zuordnungen \emph{falsch} sind, sondern weil sie von den Entscheidungsmustern des annotierenden LLM abweichen.
Die manuelle Validierung von 300~Sätzen ($\kappa = 0{,}737$ Satztyp, $\kappa = 0{,}861$ Taxonomie) zeigt zwar, dass die LLM-Annotationen grundsätzlich mit menschlichen Urteilen übereinstimmen.
Sie kann jedoch nicht ausschließen, dass in den verbleibenden 5\,317~Sätzen systematische Verzerrungen bestehen, die semantische Ansätze begünstigen.

\paragraph{Modellauswahl.}
Die Wahl des SBERT-Modells (T-Systems cross-en-de-roberta) erfolgte auf dem initialen Fulda-Datensatz, der den Großteil des finalen Datensatzes stellt.
Obwohl die Modellauswahl eine qualitative Vorauswahl unter verfügbaren deutschsprachigen Embedding-Modellen war und keine systematische Hyperparameter-Optimierung auf dem Test-Set stattfand, besteht ein geringes Risiko leicht optimistischer F1-Werte: Die Fulda-Sätze, auf denen die Modelleignung beurteilt wurde, fließen anteilig in den finalen Train-Test-Split ein.
Da der SVM-Klassifikator unabhängig auf dem finalen Trainingssplit trainiert wurde und der Test-Split ausschließlich zur Evaluation dient, ist ein direktes Leakage unwahrscheinlich, aber eine vollständige Trennung von Modellauswahl und Evaluation wäre methodisch sauberer gewesen.

\paragraph{Statistische Absicherung.}
Die paarweisen Unterschiede zwischen den drei Ansätzen wurden mit McNemar-Tests~\cite{mcnemar1947note} statistisch abgesichert (Tab.~\ref{tab:mcnemar}).
Alle sechs Vergleiche sind nach Bonferroni-Korrektur hochsignifikant ($p < 0{,}001$), sodass Stichprobenrauschen als Erklärung für die beobachteten Leistungsunterschiede ausgeschlossen werden kann.
Die Ergebnisse basieren jedoch auf Punktschätzern ohne Konfidenzintervalle; Bootstrap-Konfidenzintervalle hätten die Präzision der Schätzungen zusätzlich quantifiziert.
Bei Teilgruppen mit kleinem~$n$ (Uni Kassel: 7, TH Mittelhessen Taxonomie) sind die berichteten F1-Werte mit Vorsicht zu interpretieren.

\paragraph{Literaturgrundlage.}
Einige der zitierten Arbeiten -- insbesondere Kumar et al.~\cite{kumar2025automated}, Pangakis et al.~\cite{pangakis2023validation} und Reiss~\cite{reiss2023testing} -- sind als arXiv-Preprints ohne formales Peer-Review verfügbar.
Ihre Ergebnisse werden in dieser Arbeit als ergänzende Evidenz herangezogen, nicht als alleinige Grundlage für Entwurfsentscheidungen.

\paragraph{Taxonomie-Evaluation.}
Die Evaluation der Taxonomie-Zuordnung basiert auf der \emph{Primärstufe} (dem ersten Wert bei Mehrfachzuordnungen).
Dadurch wird die Fähigkeit des Systems, Nebenverben auf anderen Taxonomiestufen korrekt zu erkennen, nicht vollständig erfasst.
Eine Multi-Label-Evaluation würde ein differenzierteres Bild liefern.

\paragraph{Konfidenz.}
Das System zeigt Konfidenzwerte pro Satz an, die aus den SVM-Wahrscheinlichkeitsschätzungen abgeleitet werden.
Die Evaluation (Kap.~\ref{sec:eval-konfidenz}) zeigt, dass diese Werte eine sinnvolle Qualitätsindikation liefern: Korrekt klassifizierte Sätze weisen eine mittlere Typ-Konfidenz von $0{,}971$ auf, fehlklassifizierte nur $0{,}712$.
Die Werte sind jedoch nicht kalibriert: Eine angezeigte Konfidenz von 80\,\% bedeutet nicht zwangsläufig, dass 80\,\% der Sätze mit dieser Konfidenz korrekt klassifiziert werden.
Kalibrierungstechniken wie Platt-Skalierung oder isotonische Regression könnten die Aussagekraft verbessern.
